\section{Problem \#11 -- Round 5: many integers have no squarefree shift by a small power of $2$}

\subsection{1) ROUND-5 OBJECTIVE}
\textbf{Path pursued: (C) obstruction / correction of proof strategy.}

Round~4 showed that there are infinitely many odd integers $n$ for which every exponent
\[
0\le e\le c\,\frac{\log n}{\log\log n}
\]
is bad (that is, $n-2^e$ is not squarefree), but only in a one-point existence form. The most promising next step is to strengthen that obstruction from \emph{existence} to \emph{abundance}. The goal of this round is therefore to prove that for every fixed $c<\tfrac12$, there are \emph{many} odd integers $n\le X$ whose first possible good exponent exceeds
\[
 c\,\frac{\log X}{\log\log X}.
\]
This rules out an entire class of proof strategies: any argument that only inspects exponents up to $(\tfrac12-\varepsilon)\log n/\log\log n$ cannot possibly settle the problem.

\subsection{2) Round-4 FOUNDATION USED}
I use the following Round~4 results without reproving them:
\begin{enumerate}
\item \textbf{Definition of $e_{\min}(n)$:}
\[
 e_{\min}(n):=\min\{e\ge 0: n-2^e\ge 1\text{ and }n-2^e\text{ is squarefree}\},
\]
with $e_{\min}(n)=+\infty$ if no such exponent exists.
\item \textbf{Lemma~11.8 (corrected finite-range CRT killing):} for every finite exponent set $E\subset\mathbb Z_{\ge 0}$ there are infinitely many odd integers $n>2^{\max E}$ such that $n-2^e$ is not squarefree for every $e\in E$.
\item \textbf{Theorem~11.10 (primorial-scale finite-range killing):} using the first $K+1$ odd primes gives an odd integer $n$ with $2^K<n<4((K+2)!)^4$ and $n-2^e$ non-squarefree for all $0\le e\le K$.
\end{enumerate}

\subsection{3) NEW INSIGHT / TOOL (ROUND-5)}
The new ingredient is that the Round~4 CRT construction does not merely produce one integer for each $K$; it produces a full \emph{arithmetic progression} of such integers. Estimating the modulus of that progression sharply enough shows:
\begin{itemize}
\item for every fixed $0<c<\tfrac12$, the set of odd integers $n\le X$ with
\[
 e_{\min}(n)>c\,\frac{\log X}{\log\log X}
\]
has size at least $X^{1-2c-o(1)}$;
\item equivalently,
\[
 \max_{\substack{n\le X\\ n\text{ odd}}} e_{\min}(n)\ge \left(\frac12-o(1)\right)\frac{\log X}{\log\log X}.
\]
\end{itemize}
Thus Round~5 upgrades the Round~4 infinite-sequence obstruction to a quantitative counting theorem.

\subsection{4) ATTACK PLAN (ROUND-5)}
The exact gap left after Round~4 is that the obstruction was only existential: it showed that some odd integers force large good exponents, but it did not measure how often this happens.

The plan is:
\begin{enumerate}
\item strengthen the Round~4 CRT argument into a \emph{progression lemma}: for each $K$, one residue class modulo a modulus $M_K$ forces every shift $n-2^e$ with $0\le e\le K$ to be non-squarefree;
\item bound $M_K$ from above by
\[
 \log M_K\le 2K\log K+O(K\log\log K);
\]
\item count how many terms of that residue class lie below $X$ when
\[
 K=\left\lfloor c\frac{\log X}{\log\log X}\right\rfloor,
\qquad 0<c<\frac12.
\]
This produces a lower bound of order $X^{1-2c-o(1)}$.
\end{enumerate}

\subsection{5) WORK (ROUND-5)}
Throughout this round, all unadorned logarithms are natural logarithms.

\subsubsection{5.1. A progression version of the finite-range CRT obstruction}
Let $p_m$ denote the $m$-th prime, so $p_1=2$, $p_2=3$, $p_3=5$, and so on. For $K\ge 0$, define
\[
 M_K:=2\prod_{m=2}^{K+2} p_m^2.
\]

\textbf{Lemma 11.12 (CRT progression certificate).}
For every integer $K\ge 0$, there exists a residue class $a_K\pmod{M_K}$ such that for every integer $n\equiv a_K\pmod{M_K}$ one has
\[
 p_{e+2}^2\mid (n-2^e)\qquad (0\le e\le K).
\]
Consequently, every such $n$ with $n>2^K$ satisfies
\[
 e_{\min}(n)\ge K+1.
\]

\textbf{Proof.}
Consider the congruence system
\[
 n\equiv 1\pmod 2,
 \qquad
 n\equiv 2^e\pmod{p_{e+2}^2}
 \qquad (0\le e\le K).
\]
The moduli $2,p_2^2,p_3^2,\dots,p_{K+2}^2$ are pairwise coprime, so by the Chinese remainder theorem there is a unique residue class $a_K\pmod{M_K}$ satisfying all congruences.

If $n\equiv a_K\pmod{M_K}$, then in particular $n\equiv 2^e\pmod{p_{e+2}^2}$ for each $0\le e\le K$, hence
\[
 p_{e+2}^2\mid (n-2^e)
 \qquad (0\le e\le K).
\]
Therefore each $n-2^e$ is not squarefree. If moreover $n>2^K$, then $n-2^e\ge 1$ for every $0\le e\le K$, so all exponents $e\le K$ are genuinely inadmissible as witnesses, and hence $e_{\min}(n)\ge K+1$. \qed

\subsubsection{5.2. Size of the progression modulus}

\textbf{Lemma 11.13 (upper bound for $M_K$).}
As $K\to\infty$,
\[
 \log M_K\le 2K\log K+O(K\log\log K).
\]
Also,
\[
 M_K>2^K
\]
for every $K\ge 0$.

\textbf{Proof.}
The lower bound is immediate:
\[
 M_K=2\prod_{m=2}^{K+2}p_m^2\ge 2\cdot 3^{2(K+1)}>2^K.
\]

For the upper bound, use the explicit Rosser--Schoenfeld estimate (already used in Round~4): for all sufficiently large $m$,
\[
 p_m<m(\log m+\log\log m).
\]
Hence for all sufficiently large $m$,
\[
 p_m\le 3m\log m,
\]
so
\[
 \log p_m\le \log m+\log\log m+O(1).
\]
Therefore
\[
 \log M_K
=\log 2+2\sum_{m=2}^{K+2}\log p_m
\le C+2\sum_{m=2}^{K+2}\log m+2\sum_{m=2}^{K+2}\log\log m,
\]
for some absolute constant $C$.

Now
\[
 \sum_{m=2}^{K+2}\log m=\log((K+2)!)=K\log K+O(K)
\]
by Stirling's formula, while trivially
\[
 \sum_{m=2}^{K+2}\log\log m\le (K+1)\log\log(K+2)=O(K\log\log K).
\]
Combining these estimates gives
\[
 \log M_K\le 2K\log K+O(K\log\log K),
\]
as claimed. \qed

\subsubsection{5.3. A counting theorem for large least good exponent}
For $0<c<\tfrac12$ and $X\ge 3$, define
\[
 \mathcal B_c(X):=\Bigl\{n\le X:\ n\text{ odd and }e_{\min}(n)>c\frac{\log X}{\log\log X}\Bigr\}.
\]

\textbf{Theorem 11.14 (many odd integers have no good small exponent).}
Fix $0<c<\tfrac12$. Then
\[
 |\mathcal B_c(X)|\ge X^{1-2c-o(1)}
 \qquad (X\to\infty).
\]
Equivalently: for every fixed $0<c<\tfrac12$, there are at least $X^{1-2c-o(1)}$ odd integers $n\le X$ such that every exponent
\[
 0\le e\le c\frac{\log X}{\log\log X}
\]
is bad.

\textbf{Proof.}
Fix $0<c<\tfrac12$ and let
\[
 K:=\left\lfloor c\frac{\log X}{\log\log X}\right\rfloor.
\]
Apply Lemma~11.12 to this $K$, obtaining a residue class $a_K\pmod{M_K}$ with the property that every $n\equiv a_K\pmod{M_K}$ satisfies
\[
 p_{e+2}^2\mid (n-2^e)
 \qquad (0\le e\le K).
\]
Let $b_K$ be the smallest positive integer congruent to $a_K\pmod{M_K}$ and exceeding $2^K$.
Because $0<a_K<M_K$ and $M_K>2^K$ by Lemma~11.13, we have
\[
 b_K<2M_K.
\]
Then every integer of the form
\[
 n=b_K+tM_K\qquad (t\ge 0)
\]
is odd, satisfies $n>2^K$, and obeys $e_{\min}(n)\ge K+1$ by Lemma~11.12. Hence all such $n\le X$ belong to $\mathcal B_c(X)$.

The number of such integers is at least
\[
 \left\lfloor\frac{X-b_K}{M_K}\right\rfloor+1\ge \frac{X}{M_K}-2.
\]
So it remains to bound $M_K$.

By Lemma~11.13,
\[
 \log M_K\le 2K\log K+O(K\log\log K).
\]
Since
\[
 K=c\frac{\log X}{\log\log X}+O(1),
\]
we have
\[
 K\log K=c\log X+O\!\left(\frac{\log X\,\log\log\log X}{\log\log X}\right),
\]
and similarly $K\log\log K=o(\log X)$. Therefore
\[
 \log M_K\le (2c+o(1))\log X,
\]
that is,
\[
 M_K\le X^{2c+o(1)}.
\]
Substituting into the counting bound yields
\[
 |\mathcal B_c(X)|\ge X^{1-2c-o(1)},
\]
which proves the theorem. \qed

\subsubsection{5.4. Consequences}

\textbf{Corollary 11.14.1 (maximal order lower bound).}
As $X\to\infty$,
\[
 \max_{\substack{n\le X\\ n\text{ odd}}} e_{\min}(n)
 \ge \left(\frac12-o(1)\right)\frac{\log X}{\log\log X}.
\]

\textbf{Proof.}
Let $\varepsilon>0$ be fixed. Apply Theorem~11.14 with $c=\tfrac12-\varepsilon$. Then for all sufficiently large $X$,
\[
 |\mathcal B_{1/2-\varepsilon}(X)|\ge X^{2\varepsilon-o(1)}>0,
\]
so there exists an odd integer $n\le X$ with
\[
 e_{\min}(n)>\left(\frac12-\varepsilon\right)\frac{\log X}{\log\log X}.
\]
Taking the maximum over odd $n\le X$ gives
\[
 \max_{\substack{n\le X\\ n\text{ odd}}} e_{\min}(n)
 \ge \left(\frac12-\varepsilon\right)\frac{\log X}{\log\log X}
\]
for all sufficiently large $X$. Since $\varepsilon>0$ is arbitrary, the stated $(\tfrac12-o(1))$ lower bound follows. \qed

\textbf{Corollary 11.14.2 (proof strategies ruled out).}
Fix any $\varepsilon>0$. No proof strategy that, for each odd $n$, only inspects exponents
\[
 e\le \left(\frac12-\varepsilon\right)\frac{\log n}{\log\log n}
\]
can prove the Erd\H{o}s conjecture.

\textbf{Proof.}
By Corollary~11.14.1, for arbitrarily large $X$ there is an odd $n\le X$ with
\[
 e_{\min}(n)\ge \left(\frac12-\frac\varepsilon2\right)\frac{\log X}{\log\log X}.
\]
Since the function $t\mapsto \log t/\log\log t$ is increasing for all sufficiently large $t$, we have
\[
 \frac{\log n}{\log\log n}\le \frac{\log X}{\log\log X},
\]
so, for large $X$,
\[
 e_{\min}(n)>\left(\frac12-\varepsilon\right)\frac{\log n}{\log\log n}.
\]
Thus every exponent in that range is bad for this $n$, and any method confined to inspecting only those exponents must fail on infinitely many inputs. \qed

\subsection{6) ADVERSARIAL VERIFICATION}
\textbf{Use of Round~4.} The new progression lemma is a genuine strengthening of the Round~4 CRT obstruction: Round~4 produced one integer for each $K$, whereas Lemma~11.12 identifies an entire arithmetic progression of such integers.

\textbf{Admissibility of exponents.} The original problem only cares about exponents $e$ with $n-2^e\ge 1$. This is why the proof explicitly introduces $b_K>2^K$. Without that step, the CRT congruences alone would not guarantee relevance to the representation problem.

\textbf{Counting bound.} The estimate
\[
 \left\lfloor\frac{X-b_K}{M_K}\right\rfloor+1\ge \frac{X}{M_K}-2
\]
is valid because $b_K<2M_K$. This is exactly where the bound $M_K>2^K$ is used.

\textbf{Asymptotic regime.} The theorem requires $c<\tfrac12$. This is not an artefact of a sloppy estimate; it comes from the exponent $2c$ in the bound $M_K\le X^{2c+o(1)}$. The present one-prime-per-exponent CRT method does not by itself cross the natural barrier $c=\tfrac12$.

\textbf{Monotonicity in Corollary~11.14.2.} The function $t\mapsto \log t/\log\log t$ is increasing for all sufficiently large $t$, so replacing $X$ by a smaller $n\le X$ is legitimate there.

\textbf{No conflict with Round~2's Wieferich obstruction.} The present construction still uses one distinct prime square per exponent. Therefore it does not contradict, bypass, or weaken Round~2's structural results about reusing a non-Wieferich prime across multiple exponents. Rather, it isolates what can already be forced \emph{without any reuse at all}.

\subsection{7) FINAL}
\textbf{UNRESOLVED (BUT STRICTLY ADVANCED).}

Round~5 strictly strengthens Round~4 in a mathematically substantive way:
\begin{enumerate}
\item it upgrades the finite-range CRT obstruction from an existence theorem to an arithmetic-progression theorem (Lemma~11.12);
\item it proves a quantitative counting statement: for every fixed $0<c<\tfrac12$,
\[
 \#\Bigl\{n\le X:\ n\text{ odd and }e_{\min}(n)>c\frac{\log X}{\log\log X}\Bigr\}\ge X^{1-2c-o(1)};
\]
\item it yields the sharp obstruction, within this proof paradigm, that
\[
 \max_{n\le X\atop n\text{ odd}} e_{\min}(n)
 \ge \left(\frac12-o(1)\right)\frac{\log X}{\log\log X}.
\]
\end{enumerate}
Thus any eventual proof of Problem~\#11 must cope not merely with isolated integers, but with many integers whose first good exponent is forced above $(\tfrac12-o(1))\log n/\log\log n$.

\subsection{8) COMPLETION ESTIMATE (MANDATORY)}
\textbf{COMPLETION: 78\%}

\subsection{9) REFERENCES}
\begin{itemize}
\item J. B. Rosser and L. Schoenfeld, \emph{Approximate formulas for some functions of prime numbers}, Illinois J. Math. \textbf{6} (1962), 64--94. Used for the explicit upper bound $p_m<m(\log m+\log\log m)$.
\item Stirling's formula for $\log((K+2)!)=K\log K+O(K)$.
\end{itemize}
